\documentclass[11pt,a4paper]{article}
\usepackage[margin=0.85in,top=1in,bottom=1in]{geometry}
\usepackage{amsmath,amssymb,amsfonts}
\usepackage{graphicx}
\usepackage{float}
\usepackage{enumitem}
\usepackage{microtype}
\usepackage{caption}
\usepackage[hidelinks]{hyperref}
\usepackage[most,breakable]{tcolorbox}
\tcbuselibrary{breakable,skins}
\usepackage[utf8]{inputenc}
\usepackage[T1]{fontenc}
\usepackage{lmodern}
\captionsetup{font=small,labelfont=bf,skip=4pt}
\setlist{leftmargin=1.5em,topsep=2pt}

%% Breakable card — prevents content cutoff across pages
\newtcolorbox{solcard}[3]{
  breakable, enhanced,
  colback=white, colframe=gray!50, arc=2pt,
  title={\textbf{#1}\hfill\textsf{\small #3\,pts}},
  fonttitle=\normalsize\sffamily\bfseries,
  before upper={\small\textit{#2}\par\smallskip\hrule height 0.4pt\smallskip},
  top=4pt, bottom=6pt, left=7pt, right=7pt,
  parbox=false,
}

\title{\textbf{Silicon photomultiplier tube --- Solution}}
\author{\normalsize Y23 (2023) — English translation}
\date{}
\begin{document}\maketitle\vspace{-0.5em}
\begin{solcard}{A1}{For steady state, find: voltage $V_r$ at load $r$; voltage $V_0$ on each cell and parasitic capacitance $C_0$; voltage $V_1$ on the damping layers of each cell.}{0.40}

For steady state, find:


In steady state, no current flows through $r$, so $V_r = 0$. The voltage on the cells coincides with the source voltage, that is, $V_0 = V_{op}$. The capacitor $C_1$ is parallel to the resistor, so it is discharged: $V_1 = 0$.

\par\smallskip\noindent\textbf{Answer:} $$V_r = 0$$$$V_0 = V_{op}$$$$V_1 = 0$$
\end{solcard}

\begin{solcard}{A2}{Find: change in charge $\Delta{q}_0$ of capacitor $C_0$. Express your answer using $C_0$ and $\Delta{V}$; changes in charges $\Delta{q}_{1a}$ and $\Delta{q}_{2a}$ of capacitors $C_1$ and $C_2$, respectively, of each of the active cells. Express your answers using $C_1$, $C_2$, $V_{ov}$ and $\Delta{V}$; changes in charges $\Delta{q}_{1p}$ and $\Delta{q}_{2p}$ of capacitors $C_1$ and $C_2$, respectively, of each of the passive cells. Express your answers using $C_{cell}$ and $\Delta{V}$.}{0.80}

Find:


By definition of capacity: $\Delta q_0 = C_0\Delta V$. The charge change on capacitors in active cells is similar: $\Delta q_{1a} = C_1(V_{ov} + \Delta V)$ and $\Delta q_{2a} = -C_2 V_{ov}$, and in passive cells: $\Delta q_{1p} = \Delta q_{2p} = C_{cell} \Delta V$.

\par\smallskip\noindent\textbf{Answer:} $$\Delta q_0 = C_0\Delta V$$$$\Delta q_{1a} = C_1(V_{ov} + \Delta V)$$$$\Delta q_{2a} = -C_2 V_{ov}$$$$\Delta q_{1p} = \Delta q_{2p} = C_{cell} \Delta V$$
\end{solcard}

\begin{solcard}{A3}{Find the change in voltage $\Delta{V}$ across the cells and capacitor $C_0$ upon completion of the avalanche process. Express your answer in general terms through $V_{ov}$, $n$, $N$, $C_1$, $C_2$, $C_0$.}{0.70}

Let us write down the charge conservation law for positive plates $C_0$ and $C_1$. $$n\Delta q_{1a} + \Delta q_0 + (N-n)\Delta q_{1p} = 0$$$$n C_1(V_{ov} + \Delta V) + C_0\Delta V + (N-n)C_{cell} \Delta V = 0$$$$\Delta V (nC_1 + C_0 + (N-n)C_{cell}) = - nC_1 V_{ov}$$$$\Delta V = -\cfrac{ nC_1 V_{ov}}{nC_1 + C_0 + (N-n)\cfrac{C_1 C_2}{C_1 + C_2}}$$

\par\smallskip\noindent\textbf{Answer:} $$\Delta V = -\cfrac{ nC_1 V_{ov}}{nC_1 + C_0 + (N-n)\cfrac{C_1 C_2}{C_1 + C_2}}$$
\end{solcard}

\begin{solcard}{A4}{Estimate the amplitude of thermal noise across the load $r = 50~\Omega$ at room temperature if the oscilloscope bandwidth is $1~\text{GHz}=10^9~\text{Hz}$. Using the data from the problem statement and the formula you obtained from A3, find the numerical value of $\Delta{V}$ when light hits one cell.  Conclude whether the signal from one cell can be distinguished against the background of thermal noise of the load resistance.}{0.60}

Substitute $T = 300 \text{K}$ and get: $V_{noise} = 2.88\cdot10^{-5} \text{V}$. Let us express the change in voltage through the data in terms of magnitude: $$\Delta V = -\cfrac{(k+1)(1-p)\cfrac{n}{N}V_{ov}}{1 + \cfrac{kn}{N}(1-p)}$$Substitute $n = 1$ and obtain: $$\Delta V = -2.16\cdot10^{-4} \text{V}$$Note, that $\Delta V$ in magnitude noticeably larger than $V_{noise}$, so the signal can be distinguished.

\par\smallskip\noindent\textbf{Answer:} $$V_{noise} = 2.88\cdot10^{-5} \text{V}$$$$\Delta V = -2.16\cdot10^{-4} \text{V}$$$\Delta V$ is noticeably larger in absolute value than $V_{noise}$, so the signal can be distinguished.
\end{solcard}

\begin{solcard}{B1}{Find the multiplication coefficient $M$ for the TFU with the parameters from the problem statement. Express your answer in terms of $N$, $C_1$, $C_2$, $C_0$, $e$, $V_{ov}$ and give a numerical value.}{0.80}

$Q$ can be determined through the change in charges on the negative plates $C_0$ and $C_2$. $$Q = - C_0 \Delta V - (N - 1)C_{cell} \Delta V + C_2 V_{ov}$$$$M =\left( (C_0 + (N-1)C_{cell})\cfrac{ C_1 V_{ov}}{C_1 + C_0 + (N-1)C_{cell}} + C_2 V_{ov}\right)\cdot\cfrac{1}{e}$$$$M = 32.4\cdot10^3$$

\par\smallskip\noindent\textbf{Answer:} $$M =\left( \left(C_0 + (N-1)\cfrac{C_1 C_2}{C_1 + C_2}\right)\cfrac{ C_1 V_{ov}}{C_1 + C_0 + (N-1)\cfrac{C_1 C_2}{C_1 + C_2}} + C_2 V_{ov}\right)\cdot\cfrac{1}{e}$$$$M = 32.4\cdot10^3$$
\end{solcard}

\begin{solcard}{C1}{Let's consider a pulse recorded by an oscilloscope after light hits one cell. We will assume that the pulse duration is short enough to neglect the current flowing through the leakage resistance $R$, but large enough so that the avalanche process is completed and all switches $K_i$ are open. In this approximation, the circuit can be replaced by an equivalent series $RLC$ circuit. Specify its parameters.}{0.40}
\par\smallskip\noindent\textbf{Answer:} Resistance $r$, inductance $L$ and capacitance $C_{det}$
\end{solcard}

\begin{solcard}{C2}{Estimate the quality factor of the resulting $RLC$ circuit. Will oscillations occur in this circuit?}{0.30}
\par\smallskip\noindent\textbf{Answer:} Let's use the usual formula for the quality factor of an oscillatory circuit: $$Q = \cfrac{1}{R}\sqrt{\cfrac{L}{C}} = 0.14$$The resulting value is much less than 1, which means that oscillations will not occur in this circuit.
\end{solcard}

\begin{solcard}{C3}{Qualitatively depict a graph of the load voltage versus time after light hits one of the cells. Initially, the circuit was in steady state. It is not necessary to indicate characteristic values.}{0.30}
\begin{figure}[H]
  \centering
  \includegraphics[width=0.72\textwidth]{T1_S_EN_assets/fig1.jpg}
\end{figure}
\end{solcard}

\begin{solcard}{C4}{Compare the maximum absolute value of the voltage at the load $|V_{max}|$ after light hits one of the cells and the value $|\Delta{V}|$ found earlier. Which of these quantities is greater? Circle the correct ratio on your answer sheet.}{0.20}

The voltage across the resistor will always be less than the initial voltage of the capacitor.

\par\smallskip\noindent\textbf{Answer:} $$|V_{max}| < |\Delta V|$$
\end{solcard}

\begin{solcard}{D1}{Find the approximate dependence of the voltage on the cell's damping layer on time after light hits it (under initial steady-state conditions). Take advantage of the fact that in the conditions of this problem $\Delta{V}\ll{V_{ov}}$, $R\gg{r}$. Consider that a lot of time has passed and the voltage on the external load is almost zero. Express your answer in terms of $V_{ov}$, $R$, $C_1$, $C_2$, $t$.}{1.00}

Due to the data relationships in the condition, the voltage at $r$ and $L$ can be neglected. That is, for the restoration process the cell is equivalent to an RC circuit.


Initially the voltage at $C_1$ is $V_{ov}$.  Then time dependence: $$V_1 = V_{ov}\cdot e^{-\cfrac{t}{R(C_1 + C_2)}}$$

\par\smallskip\noindent\textbf{Answer:} $$V_1 = V_{ov}\cdot e^{-\cfrac{t}{R(C_1 + C_2)}}$$
\begin{figure}[H]
  \centering
  \includegraphics[width=0.72\textwidth]{T1_S_EN_assets/fig2.jpg}
\end{figure}
\end{solcard}

\begin{solcard}{D2}{Find the change in voltage $\Delta{V}$ across capacitor $C_0$ when light hits one ($n=1$) of the unrecovered cells, the voltage on the quenching layer of which was $V_1$.}{0.50}

The answer can be obtained in a similar way to part A, but now the change in voltage by $C_2$ will be $V_{ov} - V_1$ instead of just $V_{ov}$.

\par\smallskip\noindent\textbf{Answer:} $$\Delta V = -\cfrac{ C_1( V_{ov} - V_1)}{C_1 + C_0 + (N-1)\cfrac{C_1 C_2}{C_1 + C_2}}$$
\end{solcard}

\begin{solcard}{E1}{Let $p(t)$ be the probability that during time $t$ the cell did not fire even once. Find $p(t + dt)$, express the answer in terms of $p(t)$, $dt$, $T$. Consider that $dt \ll T$.}{0.30}

The probability that the cell will not work in time $dt$ is $1 - \cfrac{dt}{T}$, therefore: $$p(t+dt) = p(t) \cdot\left(1 - \cfrac{dt}{T}\right)$$

\par\smallskip\noindent\textbf{Answer:} $$p(t+dt) = p(t) \cdot\left(1 - \cfrac{dt}{T}\right)$$
\end{solcard}

\begin{solcard}{E2}{Find the probability that the cell will not fire even once during time $t$. Express your answer using $t$, $T$.}{0.20}

Let's use the result of the previous paragraph. $$dp = -p\cdot\cfrac{dt}{T}$$$$\cfrac{dp}{p} = -\cfrac{dt}{T}$$Use $p(0) = 1$, we get:

\par\smallskip\noindent\textbf{Answer:} $$p(t) = e^{-\cfrac{t}{T}}$$
\end{solcard}

\begin{solcard}{E3}{Let the cell fire at time $t = 0$. Find the probability that the next time it will work in the time interval from $t$ to $t + dt$.}{0.20}

The desired expression is the product of the probability that the cell will not work in time $t$, and the probability of operation in $dt$.

\par\smallskip\noindent\textbf{Answer:} $$p(t, t+dt) = e^{-\cfrac{t}{T}} \cdot \cfrac{dt}{T}$$
\end{solcard}

\begin{solcard}{E4}{Find the average amplitude of the photoresponse pulses $\mu$. Consider that different cells produce impulses independently of each other.}{0.40}

From the definition of average: $$\mu = \int\limits_0^{+\infty} V(t)\cdot p(t) \cdot \cfrac{dt}{T} = \int\limits_0^{+\infty} \left(1 - e^{-\cfrac{t}{\tau}}\right)\cdot e^{-\cfrac{t}{T}} \cdot \cfrac{dt}{T} = \int\limits_0^1 \left(1 - x^{\cfrac{T}{\tau}}\right)dx$$Where $x = e^{-\cfrac{t}{T}}$. $$\mu = 1 - \cfrac{1}{\cfrac{T}{\tau} +1} = \cfrac{T}{T + \tau}$$

\par\smallskip\noindent\textbf{Answer:} $$\mu = \cfrac{T}{T + \tau}$$
\end{solcard}

\begin{solcard}{E5}{Find the standard deviation of the photoresponse pulse amplitude $\sigma$. Express your answer in terms of $\tau$, $T$. Note: The standard deviation can be calculated using the formula $\sigma=\sqrt{\mu(V^2)-\mu^2(V)}$, where $\mu(V^2)$ is the average amplitude squared, $\mu^2(V)$ is the average amplitude squared.}{0.60}

Note: The standard deviation can be calculated using the formula $\sigma=\sqrt{\mu(V^2)-\mu^2(V)}$, where $\mu(V^2)$ is the average amplitude squared, $\mu^2(V)$ is the average amplitude squared.


$\mu(V)$ was found in the previous point, similarly we will find $\mu(V^2)$: $$\mu(V^2) = \int\limits_0^1\left(1 - x^{\cfrac{T}{\tau}}\right)^2dx = 1 - \cfrac{2}{\cfrac{T}{\tau}+1} + \cfrac{1}{\cfrac{2T}{\tau}+1}$$$$\mu(V^2) - \mu(V)^2 = \cfrac{1}{\cfrac{2T}{\tau}+1} - \cfrac{1}{\left(\cfrac{T}{\tau}+1\right)^2} = \cfrac{\cfrac{T^2}{\tau^2}}{\left(\cfrac{T}{\tau}+1\right)^2\left(\cfrac{2T}{\tau}+1\right)} = \left(\cfrac{T}{T+\tau}\right)^2\cdot\cfrac{\tau}{2T+\tau}$$$$\sigma = \cfrac{T}{T+\tau}\cdot\sqrt{\cfrac{\tau}{2T + \tau}}$$

\par\smallskip\noindent\textbf{Answer:} $$\sigma = \cfrac{T}{T+\tau}\cdot\sqrt{\cfrac{\tau}{2T + \tau}}$$
\end{solcard}

\begin{solcard}{E6}{Find $SNR$ amplitudes of photoresponse pulses under the influence of uniform background illumination in the following cases: $\tau=0$; $\tau\gg{T}$; $\tau=T$.}{0.30}

Find $SNR$ amplitudes of photoresponse pulses under the influence of uniform background illumination in the following cases:


$$SNR = \sqrt{\cfrac{2T + \tau}{\tau}}$$

\par\smallskip\noindent\textbf{Answer:} $$SNR(\tau = 0) = \infty$$$$SNR(\tau \gg T) = 1$$$$SNR(\tau = T) = \sqrt{3}$$
\end{solcard}

\end{document}